\section{Constants and uniformity}\label{app:constants}
Every estimate in the negative proof is uniform over the full affine
space of permitted Gram corrections after positivity has been imposed.
The following table records each bound and the only inputs it depends
on. The dimension and rank of a Gram factor are never among them, and
no exponential is evaluated at an astronomical argument: the tail bounds
\eqref{eq:dyadic-tail-budget} need only $2^{k-2}\ge8k+40$ for $k\ge16$.

\begin{longtable}{@{}>{\raggedright\arraybackslash}p{.28\textwidth}>{\raggedright\arraybackslash}p{.67\textwidth}@{}}
\toprule
Estimate & Bound and dependence\\
\midrule\endhead
Fixed test & $1024$ points, $r\le2$, $j\le256$;
$\|c\|_1<2^{30}$ and $c^*S_Fc<-2^{42}$.
Only the integer formula \eqref{eq:witness}.\\[5pt]
Fixed coefficients & $16(p+1)(q+1)(r+1)(s+1)$ for the expanded
coefficients of $Q_{0,\infty}$ and $L_\infty$.
Only the explicit baseline matrices.\\[5pt]
Mixed correction & $4\sqrt{(p+1)(q+1)(r+1)(s+1)}$.
Only the fixed diagonal and positive semidefiniteness.\\[5pt]
Truncation & $2^{30}\epsilon^{-8}e^{-\epsilon m/4}$.
Only $m,\epsilon$, with $0<\epsilon\le1$.\\[5pt]
Repeated-node correction & $2^{26}\epsilon^{-8}e^{-\epsilon m/2}$.
The same parameters and the complete low means.\\[5pt]
Combined tail & $\mathcal E=2^{32}\epsilon^{-8}e^{-\epsilon m/4}$.
For $m=\epsilon^{-2}$, $\epsilon=2^{-k}$, $k\ge16$:
$\mathcal E\le\epsilon^{-2}$ and $\epsilon^4\mathcal E\le\epsilon$.\\[5pt]
Boundary bounds & $2^{20}\epsilon^{-4}$ globally and
$2^{20}\epsilon^{-2}$ on the real rectangle.
The latter uses the sharper $17\epsilon^{-2}$ estimate and
$\epsilon\le2^{-12}$.\\[5pt]
Rectangle weights & $\theta_1=2^{-13}$, $\theta_2=2^{-5}$;
$2\theta_1\theta_2=2^{-17}$. Only the fixed rectangles.\\[5pt]
Pure-kernel improvement & $2^{20}\epsilon^{2^{-17}}$
after multiplication by $\epsilon^4$. Only the preceding two bounds.\\[5pt]
Baseline comparison & $2^{20}\epsilon$ before truncation,
$2^{30}\epsilon$ after truncation. Only the fixed speeds, $\epsilon\le2^{-13}$, and $k\ge16$.\\[5pt]
Scalar losses & $2^{90}\epsilon$ and $2^{80}\epsilon^{2^{-17}}$.
The entrywise bounds multiplied by $\|c\|_1^2<2^{60}$.\\
\bottomrule
\end{longtable}

The passage from $\epsilon$ to $m$ accounts for the exponent in
Theorem~\ref{thm:finite-scale}:
\[
 \epsilon=m^{-1/2},\qquad
 \epsilon^{2^{-17}}=m^{-1/262144}.
\]
At $k=38\cdot2^{17}+1=\ScaleExponent$, the slow error equals
$2^{42}2^{-2^{-17}}\le2^{42}-2^{24}$, by the convexity inequality
used in Corollary~\ref{cor:threshold}. The other error is less than one,
so the budget is at most $-2^{24}+1<0$. An entirely integer check of
the slow-error bound is
\[
 (2^{18}-1)^{2^{17}}>2^{18\cdot2^{17}-1}.
\]
At the preceding integer $k=38\cdot2^{17}$ the slow error is exactly
$2^{42}$ and the budget is positive. Both errors decrease with $k$,
so $\ScaleExponent$ is the least integer making this particular budget
negative.
